\clearpage
\section{Заключение}
Исследование показывает, как внутренний индекс меняет выбор разбиения при ограниченной внешней информации. На четырёх из шести наборов максимум silhouette совпадает с лучшим проверочным кандидатом. В этих условиях внешний допуск иногда препятствует его выбору. На AG News и Dry Bean использование внешних меток повышает качество относительно геометрического максимума. При этом на Dry Bean адаптивное правило уступает внешнему лидеру и выбору по бутстрепным победам. Преимущество одного способа сочетания оценок на всех исследованных данных не установлено.

Основной результат анализа --- разделение потери на исключение качественных кандидатов и выбор внутри допуска. На Digits и 20 Newsgroups при 20 метках действует первый механизм, на Dry Bean преобладает второй, на AG News важны оба. Это объясняет, почему расширение допуска может как помочь, так и ухудшить выбор. Детерминированные границы связывают потерю с ошибкой парного сравнения, а частота и величина переключений раскрывают структуру среднего эффекта.

Парный синтетический опыт связывает изменение эффекта со смыслом целевых классов при одинаковых признаках. Проверка состава разметки обнаруживает различия при одном бюджете, однако покрытие меняется совместно с числом пар и представленными категориями. Повторение бутстрепа выявляет чувствительность решения около границы допуска даже при большом числе перевыборок. Поэтому оценка размера разметки дополняется анализом её состава и вычислительной устойчивости.

Полученные выводы представляют ретроспективную диагностику фиксированных задач. Практическое определение подходящих условий до раскрытия проверочных меток требует отдельного исследования. Следующие проверки должны включать независимые данные, несколько разделений, разные составы кандидатов и настройку параметров на отдельной серии. Для выяснения самостоятельной пользы адаптивного допуска важно сопоставить его с фиксированным порогом при близкой частоте переключений.
